El contraste entre el análisis asintótico formal y la evaluación experimental realizada demuestra que la complejidad teórica describe la tasa de crecimiento ante instancias masivas, pero son las constantes ocultas y la gestión de memoria las que determinan el rendimiento práctico en escalas moderadas. En el ordenamiento, la robustez de \textit{std::sort} y la estabilidad lograda en \textit{Quick Sort} mediante la partición de Hoare y pivote central demuestran cómo optimizaciones a nivel de punteros y reducción de llamadas en la pila superan considerablemente a algoritmos insensibles a la entrada como \textit{Merge Sort}, el cual depende $\Theta(N)$ en memoria auxiliar y termina siendo menos eficaz. A su vez, la asimetría observada en \textit{Patience Sort} ratifica que el rendimiento real puede oscilar entre un tiempo casi lineal $\mathcal{O}(N)$ y una degradación severa en función de la estructura y orden de los datos. En el ámbito matricial, el método de \textit{Strassen} confirmó empíricamente su menor orden asintótico frente al crecimiento cúbico $\Theta(N^3)$ del enfoque tradicional, pero exhibió un umbral de corte claro ($N \ge 256$): por debajo de este tamaño, el costo constante de particionar bloques y alojar submatrices temporales supera el beneficio algebraico de ahorrar multiplicaciones escalares. En consecuencia, la selección de un algoritmo óptimo exige sopesar no solo su cota asintótica, sino que tambien cualidades de la entrada y todo el contexto en el que se quiere aplicar un algoritmo como solución.